\section{Agent Frameworks and Infrastructure}
\label{sec:agent-frameworks-infrastructure}

Agent frameworks are the practical substrate on which contemporary AI self-evolution systems are built.  They define how language models are wrapped as agents, how agents invoke tools, how memory is read and written, how execution is monitored, how multiple agents coordinate, and how failures are surfaced to either another agent or a human operator.  In early discussions of self-evolving AI it is tempting to focus only on the optimization algorithm: evolutionary search, reinforcement learning, self-reflection, program synthesis, or benchmark-driven prompt improvement.  Yet the algorithm alone does not make a self-evolving system.  A system can only improve itself if it can repeatedly run, observe its own behavior, preserve evidence, compare variants, prevent regressions, and deploy improved configurations without destroying the context that made the improvement meaningful.  These capabilities are provided, or at least constrained, by the framework layer.

This chapter surveys eight representative open-source agent and LLM-programming frameworks: AutoGPT, MetaGPT, AutoGen, CrewAI, DSPy, CAMEL, LangGraph, and OpenHands.  They were selected because they cover the main design poles in the current ecosystem.  AutoGPT represents the early autonomous-agent loop and the transition from viral prototype to platform.  MetaGPT represents role-based standard operating procedures for multi-agent software production.  AutoGen represents a message-driven, conversational, actor-like runtime for multi-agent systems.  CrewAI represents a lightweight production-oriented split between autonomous crews and deterministic flows.  DSPy represents declarative LLM programming in which prompts and demonstrations become optimizable program parameters.  CAMEL represents role-playing societies, critic-in-the-loop collaboration, and research-oriented agent environments.  LangGraph represents graph-structured, stateful, cyclic workflows with checkpointing.  OpenHands represents the full software-development environment: an agent controlling an editor, terminal, browser, and sandbox.

The unifying observation is that none of these frameworks is, by itself, a complete self-evolution layer.  Each provides one or more ingredients.  Some provide orchestration, some memory, some tool execution, some evaluation, and some optimization.  But current frameworks generally lack persistent lineage tracking, long-horizon cross-run memory, regression-safe promotion, unified experiment provenance, and benchmark governance.  Therefore, a self-evolving AI architecture should not replace these frameworks.  It should sit above them as an evolution layer that treats framework instances as execution substrates, collects traces and metrics, proposes mutations, evaluates variants, and promotes only those changes whose evidence survives regression checks.  In this sense, agent frameworks are to self-evolution what operating systems and build systems are to software engineering: they are necessary infrastructure, but they are not the entire engineering discipline.

% Auto-generated by scripts/analyze_github_project_data.mjs
\subsection{GitHub Project Corpus and Sampling Funnel}
\label{subsec:github-project-corpus-sampling}
The GitHub corpus is organized as a funnel rather than a flat awesome list. The raw discovery layer contains 416 timestamp-indexed GitHub captures. The classification layer contains 416 repositories with category, theme, stack, and time-slice annotations. The paper-facing analysis layer currently contains 129 repositories that receive public project reports and model-card pages. Within the classified corpus, 71 repositories are strict evolution-theme repositories, while 166 match a broader evolution-related keyword set including self-improvement, reflection, recursive improvement, prompt evolution, and code evolution.

\begin{table}[t]
\centering
\caption{GitHub project corpus funnel.}
\label{tab:github-project-corpus-funnel}
\begin{tabular}{lrp{0.55\linewidth}}
\toprule
Layer & Count & Meaning \\
\midrule
Raw captures & 416 & Timestamp-indexed records from \texttt{raw-github/}. \\
Classified repositories & 416 & Rows with category, theme, stack, evidence, and time-slice fields. \\
Analyzed model-card projects & 129 & Repositories promoted into site and paper-facing project analysis. \\
Strict evolution theme & 71 & Classified rows whose base theme is \texttt{evolution}. \\
Broad evolution-related & 166 & Rows matching evolution, self-improvement, reflection, or recursive-improvement signals. \\
\bottomrule
\end{tabular}
\end{table}

\begin{table}[t]
\centering
\caption{Raw GitHub category and theme distribution.}
\label{tab:github-category-theme-distribution}
\begin{tabular}{lrlr}
\toprule
Category & Count & Theme & Count \\
\midrule
框架/framework & 118 & evaluation & 89 \\
评测/evaluation & 96 & memory & 77 \\
教程/tutorial & 75 & evolution & 71 \\
工具/tool & 56 & framework & 43 \\
应用/application & 47 & education-list & 35 \\
论文代码/paper-code & 24 & research-agent & 31 \\
 &  & prompt-optimization & 26 \\
 &  & skill & 20 \\
 &  & coding-agent & 17 \\
 &  & workflow-automation & 6 \\
 &  & safety & 1 \\
\bottomrule
\end{tabular}
\end{table}

The time analysis deliberately separates repository creation from activity timestamps. The raw corpus time-slice field is an activity/content timestamp extracted from GitHub captures, so it measures when the captured page exposed recent activity. The analyzed-project timeline uses GitHub API \texttt{created\_at} where available. When GitHub API creation time is unavailable, the table marks a weaker \texttt{local\_git\_first\_commit} fallback; this is a first-observed mirror timestamp, not a release-date claim. In the current run, 49 of 129 analyzed projects have either verified creation time or a local fallback.

\begin{table*}[t]
\centering
\small
\caption{First 24 analyzed projects by verified creation month or local first-observed fallback.}
\label{tab:analyzed-project-release-timeline}
\begin{tabular}{lllp{0.21\linewidth}p{0.30\linewidth}}
\toprule
Month & Source & Repository & Category & Evolution pattern \\
\midrule
2022-09 & github\_api & \texttt{carperai/openelm} & 进化式 Prompt 优化 & 进化/搜索循环 → 评估器/打分器 \\
2023-01 & github\_api & \texttt{stanfordnlp/dspy} & 声明式 Prompt 优化 & 反馈-精炼 → 进化/搜索循环 → 评估器/打分器 \\
2023-03 & github\_api & \texttt{Significant-Gravitas/AutoGPT} & 自主 Agent 平台 & 智能体编排 \\
2023-03 & github\_api & \texttt{camel-ai/camel} & 角色扮演 Agent 框架 & 智能体编排 → 反馈-精炼 \\
2023-03 & github\_api & \texttt{noahshinn/reflexion} & 反思记忆 & 进化/搜索循环 → 反思记忆 → 反馈-精炼 → 评估器/打分器 → 训练/数据循环 \\
2023-03 & github\_api & \texttt{madaan/self-refine} & 反馈精炼 & 反馈-精炼 \\
2023-04 & local\_git\_first\_commit & \texttt{automl/auto-sklearn} & AutoML 框架 & 进化/搜索循环 → 评估器/打分器 \\
2023-04 & github\_api & \texttt{microsoft/CoML} & ML 知识库驱动 & 反馈-精炼 → 评估器/打分器 \\
2023-06 & github\_api & \texttt{FoundationAgents/MetaGPT} & 多 Agent 协作框架 & 智能体编排 → 反馈-精炼 \\
2023-07 & github\_api & \texttt{aiwaves-cn/agents} & 数据驱动 Agent 进化 & 进化/搜索循环 → 评估器/打分器 → 智能体编排 \\
2023-08 & github\_api & \texttt{langchain-ai/langgraph} & 图式 Agent 编排 & 智能体编排 \\
2023-08 & github\_api & \texttt{microsoft/autogen} & 多 Agent 对话框架 & 智能体编排 \\
2023-10 & github\_api & \texttt{google-deepmind/opro} & LLM 作为优化器 & 进化/搜索循环 → 评估器/打分器 \\
2023-10 & github\_api & \texttt{crewAIInc/crewAI} & 多 Agent 协作框架 & 智能体编排 \\
2023-11 & github\_api & \texttt{google-deepmind/funsearch} & 进化式数学发现 & 进化/搜索循环 → 评估器/打分器 \\
2024-03 & github\_api & \texttt{All-Hands-AI/OpenHands} & AI 软件开发平台 & 智能体编排 \\
2024-04 & github\_api & \texttt{princeton-nlp/SWE-agent} & 软件工程 Agent & 反馈-精炼 → 评估器/打分器 \\
2024-07 & github\_api & \texttt{shengranhu/adas} & Agent 架构自动搜索 & 进化/搜索循环 → 智能体编排 → 评估器/打分器 \\
2024-08 & local\_git\_first\_commit & \texttt{alfa-group/tutorial\_gp\_llm} & GP+LLM 教学 & 进化/搜索循环 → 评估器/打分器 \\
2024-09 & local\_git\_first\_commit & \texttt{OpenBMB/AgentVerse} & 多 Agent 仿真平台 & 智能体编排 → 反思记忆 \\
2024-10 & local\_git\_first\_commit & \texttt{siyuyuan/evoagent} & 进化式多 Agent 系统 & 进化/搜索循环 → 智能体编排 \\
2024-11 & local\_git\_first\_commit & \texttt{xiaofangxd/LLM\_EA} & LLM+EA 交叉综述 & 文献综述 \\
2025-03 & local\_git\_first\_commit & \texttt{wuxingyu-ai/LLM4EC} & LLM+EC 交叉综述 & 文献综述 \\
2025-05 & github\_api & \texttt{DeepAuto-AI/automl-agent} & 多 Agent AutoML & 智能体编排 → 进化/搜索循环 → 评估器/打分器 \\
\bottomrule
\end{tabular}
\end{table*}

\begin{table}[t]
\centering
\caption{Top strict evolution-theme repositories by local star snapshot.}
\label{tab:top-evolution-theme-repos}
\begin{tabular}{lrrp{0.42\linewidth}}
\toprule
Repository & Stars & Time & Description \\
\midrule
\texttt{aiming-lab/AutoResearchClaw} & 12600 & 2026-05 & AutoResearchClaw 是从 research idea 到 paper 的自主/协作式科研 agent 管线，结合多 agent debate、实验沙箱、claim verification、HITL co-pilot、MetaClaw cross-run learning 和 ARC-Bench。 \\
\texttt{aden-hive/hive} & 10400 & 2026-05 & Aden Hive 是生产 AI 的 multi-agent harness，强调状态持久化、崩溃恢复、成本控制、审计轨迹、MCP 工具集成和失败驱动的 graph evolution。 \\
\texttt{metauto-ai/gptswarm} & 998 & 2026-05 &  The First Self-Improving agents with RL / Prompting Optimization \\
\texttt{adolfousier/opencrabs} & 755 & 2026-05 & OpenCrabs 是受 OpenClaw 启发的单二进制多渠道 AI agent，强调本地 brain files、memory search、custom commands、cron/background jobs 和 self-update 组成的自改进循环。 \\
\texttt{human-agent-society/coral} & 667 & 2026-05 & CORAL 是面向 open-ended discovery 的多代理自主演化基础设施，用隔离 git worktrees、共享状态目录、grader daemon、heartbeat prompt 和多 runtime 集成推动 agent teams 连续改进。 \\
\texttt{a-evo-lab/a-evolve} & 552 & 2026-05 & A-Evolve 是通用 self-improving agent 基础设施：给定 base agent、benchmark 和 evolution algorithm，就把 prompt、skills、memory 等 agent workspace 文件作为可变状态进行迭代。 \\
\texttt{thudm/webrl} & 524 & 2026-05 & Building Open LLM Web Agents with Self-Evolving Online Curriculum RL \\
\texttt{ecnu-icalk/autoskill} & 424 & 2026-05 & AutoSkill: Experience-Driven Lifelong Learning via Skill Self-Evolution \\
\texttt{china-qijizhifeng/agentic-Harness-engineering} & 391 & 2026-05 & Agentic Harness Engineering treats prompts, tools, middleware, skills, sub-agents, memory, and evaluator scaffolds as evolvable harness surfaces while the base model stays fixed. \\
\texttt{feiliu36/eoh} & 319 & unknown & Evolution of Heuristics \\
\texttt{meituan/EvoCUA} & 317 & 2026-05 & 美团 EvoCUA 计算机使用 Agent 项目，发布 EvoCUA-32B/8B 并在 OSWorld、WindowsAgentArena 等 GUI 自动化评测上报告结果。 \\
\texttt{brain-research/guided-evolutionary-strategies} & 273 & unknown & Guided Evolutionary Strategies \\
\bottomrule
\end{tabular}
\end{table}


\subsection{Framework Overview}
\label{subsec:framework-overview}

The surveyed frameworks differ along six practical dimensions.  The first is the orchestration pattern: whether the framework centers on a loop, a graph, a conversation, a role-based pipeline, or a declarative program.  The second is memory: whether the framework treats memory as a simple chat history, a vector store, a persistent checkpoint, an entity graph, or an explicit optimization artifact.  The third is tool use: whether tools are command functions, external workbenches, code interpreters, browser controllers, terminal actions, or full sandbox operations.  The fourth is evaluation: whether the framework has a built-in benchmark layer, relies on external tests, or treats evaluation as a user-supplied metric.  The fifth is extensibility: whether users extend the system by writing roles, blocks, graph nodes, signatures, skills, tools, or complete agents.  The sixth is support for self-evolution: whether the framework can naturally express repeated variation, measurement, memory, and promotion.

\begin{table*}[t]
\centering
\small
\caption{High-level comparison of representative agent frameworks and infrastructure for AI self-evolution.}
\label{tab:agent-framework-comparison}
\begin{tabular}{p{0.11\textwidth}p{0.15\textwidth}p{0.15\textwidth}p{0.15\textwidth}p{0.13\textwidth}p{0.14\textwidth}p{0.12\textwidth}}
\hline
\textbf{Framework} & \textbf{Orchestration pattern} & \textbf{Memory model} & \textbf{Tool use} & \textbf{Evaluation support} & \textbf{Extensibility model} & \textbf{Self-evolution support} \\
\hline
AutoGPT & Think--Act--Observe loop; modern platform adds block workflows and agent runtime. & Short-term context, local cache, embedding retrieval, platform state. & Commands, web search, files, code execution, platform blocks. & Agent Protocol and benchmark integrations; platform monitoring. & Commands, blocks, Forge components, platform agents. & Good for autonomous trial loops, weak for disciplined lineage and regression governance without extra infrastructure. \\
MetaGPT & Role--Action--Message pipeline based on software-company SOPs. & Multi-layer role memory, long-term memory modules, skill loading, document artifacts. & Role-specific actions, search, code generation, project artifacts, experimental extensions. & Project-level output checks, agbenchmark use, SELA and AFlow evaluation modules. & Roles, actions, prompts, SOPs, extensions such as SELA. & Strong conceptual fit for evolving team procedures; needs persistent experiment ledger and promotion policy. \\
AutoGen & Actor-like message runtime plus high-level conversational teams. & Memory and model-context abstractions; list memory and conversation state. & Function tools, workbenches, tool agents, code executors, extensions. & agbench and task-specific evaluations; Studio for inspection. & Core agents, AgentChat teams, tools, handoffs, extensions. & Strong for multi-agent deliberation and variant conversations; needs cross-run selection and regression prevention. \\
CrewAI & Crew + Flow dual architecture: autonomous teams plus event-driven production flows. & Short-term, long-term, entity, and knowledge memory. & Agent tools, custom decorators, task-specific tool assignment, flow steps. & User-defined task outputs, control-plane observability, external tests. & Agents, tasks, crews, processes, flows, tools. & Strong for packaging recurring workflows; evaluation and lineage remain external concerns. \\
DSPy & Declarative Signature--Module--Optimizer programming. & Examples, traces, demonstrations, optimizer state, LM call history. & Predictors can call tools via modules; focus is program optimization. & Central: metrics, dev sets, Evaluate, optimizer compile loops. & Signatures, modules, adapters, teleprompters, optimizers. & Strongest built-in optimization loop for prompts and modules; weaker as a full autonomous runtime. \\
CAMEL & Role-playing societies and workforce coordination. & Chat-history memory, score-based context selection, memory records, multiple storage backends. & Toolkits, interpreters, browser and terminal tools, MCP agents, verifiers. & Critic agents, environments, verifiers, task outcomes. & Agents, societies, workforces, environments, toolkits, storage backends. & Strong for critic-in-the-loop and synthetic data; needs systematic variant tracking and promotion gates. \\
LangGraph & Directed state graph with cycles, branches, channels, and Pregel-style execution. & State channels and checkpoints with SQLite/PostgreSQL persistence. & LangChain-compatible tools, prebuilt ReAct agents, graph nodes. & External metrics; Studio/Cloud aid inspection and replay. & Nodes, edges, conditional routing, channels, checkpoint stores. & Excellent substrate for iterative refinement loops; evolution policy must be layered above graph execution. \\
OpenHands & Controller--Agent--Environment loop for software engineering tasks. & Conversation and workspace state, microagent/skill context, enterprise storage. & Editor, file system, terminal, browser, Docker sandbox, integrations. & Software-task benchmarks and CI-style verification are natural but not automatic. & Skills, microagents, runtime integrations, enterprise extensions. & Strong environment for code self-improvement; requires external evaluator, lineage, and rollback system. \\
\hline
\end{tabular}
\end{table*}

The table shows that the field has converged on a set of recurring primitives.  Nearly every framework needs an agent abstraction, a message or state abstraction, a tool abstraction, and an execution boundary.  What differs is the first-class object.  In AutoGPT, the first-class object was originally an autonomous loop pursuing a goal.  In MetaGPT, it is the role within a standard operating procedure.  In AutoGen, it is the message-addressable agent.  In CrewAI, it is the crew or flow.  In DSPy, it is the declarative module compiled against a metric.  In CAMEL, it is the social role in a structured dialogue.  In LangGraph, it is the graph state and its transitions.  In OpenHands, it is the agent acting inside a real development environment.

For self-evolution, the crucial question is not which first-class object is aesthetically preferable, but which object can be mutated, evaluated, and promoted safely.  AutoGPT can mutate goals, commands, and block graphs.  MetaGPT can mutate role definitions, action prompts, and SOP transitions.  AutoGen can mutate team composition, speaker selection, handoff rules, and tool workbenches.  CrewAI can mutate tasks, process modes, flows, and tool assignments.  DSPy can mutate instructions, demonstrations, signatures, modules, and optimizer choices.  CAMEL can mutate role prompts, critic policies, planning structures, and generated data.  LangGraph can mutate nodes, edges, routing predicates, checkpoint strategies, and retry policies.  OpenHands can mutate skills, microagents, code patches, and environment procedures.  All of these mutations are meaningful.  None is safe unless accompanied by evaluation records, lineage identifiers, and rollback rules.

A second cross-cutting pattern is the distinction between \emph{in-run adaptation} and \emph{cross-run evolution}.  Many frameworks support in-run adaptation: an agent sees an observation and changes its next action; a graph routes to a repair node; a critic asks for revision; a DSPy optimizer selects better demonstrations during compilation.  Self-evolution, however, requires cross-run persistence.  The system must remember that variant A beat variant B on benchmark C under seed D with model E, tool version F, and cost G.  It must know whether the improvement held under a rerun, whether it regressed another task, and whether it should become the new baseline.  Present frameworks often store enough state to continue a single task, but not enough evidence to govern a lineage of improving systems.

A third pattern is that agent frameworks expose different levels of control.  AutoGPT and OpenHands expose rich environment control, including potentially dangerous file, shell, and web actions.  DSPy exposes optimization control but less environment control.  LangGraph exposes deterministic workflow control but delegates many details to nodes.  AutoGen and CAMEL expose conversation control.  MetaGPT and CrewAI expose organizational control.  A mature self-evolution platform will likely combine several of these layers: DSPy-style metric compilation for prompts, LangGraph-style cyclic workflows for evaluation pipelines, OpenHands-style sandboxed software execution, AutoGen-style deliberation among reviewers, and MetaGPT- or CrewAI-style role specialization.

\subsection{AutoGPT: From TAO Loop to Hype-Aware Platform}
\label{subsec:autogpt}

AutoGPT is historically important because it popularized the autonomous-agent pattern for the broader developer community.  Its classic design can be summarized as a Think--Act--Observe loop.  The user states a goal.  The model thinks about the goal and current context.  It selects an action such as searching the web, writing a file, running code, or reading memory.  The environment returns an observation.  The observation is appended to context or memory, and the loop repeats.  This pattern is simple enough to explain in a diagram and powerful enough to produce the early impression of autonomous behavior.  It became the default mental model for many later agent systems.

The TAO loop matters for self-evolution because it establishes a minimal closed feedback cycle.  Any evolving agent needs a loop in which hypotheses are generated, interventions are made, outcomes are observed, and subsequent behavior changes.  AutoGPT's classic loop provides this skeleton.  It does not, however, provide the full scientific apparatus needed for evolution.  The early loop tends to optimize locally for task progress, not globally for durable capability improvement.  It often lacks rigorous baselines, controlled variants, statistical reruns, and explicit lineage.  The agent may learn in the weak sense that its immediate context changes, but it does not necessarily learn in the strong sense that the system produces a persistent, validated successor.

The project materials record AutoGPT as having roughly 175,000 GitHub stars, placing it among the most visible open-source AI projects of the agent boom.  Such visibility is both a strength and a warning signal.  It demonstrates that the project defined a category and attracted enormous attention.  It also means that star count is not a clean measure of engineering health.  The reported stars-per-contributor ratio of about 213 is suspiciously high compared with more workmanlike infrastructure projects.  A high ratio can indicate broad public curiosity, viral media exposure, or speculative enthusiasm rather than proportional depth of maintenance.  In survey terms, AutoGPT should be credited for agenda-setting influence, but not treated as automatically superior because of stars alone.

The difference between hype and contribution is visible in the project's evolution.  The original AutoGPT was a proof-of-concept autonomous agent with a command system, memory, and workspace.  The later AutoGPT platform is more ambitious: a Docker-based runtime, a Python backend, a React frontend, a market for reusable agents, analytics, monitoring, and block-based workflow construction.  The project therefore moved from a single-agent CLI imagination to a platform architecture.  This shift is important.  It shows that real users needed reliability, deployment, monitoring, sharing, and control planes more than they needed an endlessly looping agent.  The lesson for self-evolution is that the path from impressive demo to durable system runs through infrastructure.

AutoGPT's orchestration remains conceptually loop-centered.  Even when block workflows are introduced, the system is still organized around agents pursuing goals through a sequence of actions.  This makes it well suited for exploratory tasks where the path is unknown.  It is less well suited for tasks requiring strict reproducibility unless additional constraints are imposed.  A self-evolution controller using AutoGPT as a substrate would need to bound the loop, define allowable mutations, capture every action and observation, and separate exploration runs from promotion runs.  Otherwise, the same autonomy that makes AutoGPT exciting can make its improvements impossible to audit.

Its memory design is likewise useful but incomplete for evolution.  Classic AutoGPT uses short-term context and local cache or embedding-based retrieval.  This helps the agent continue a task and recall relevant prior facts.  But evolutionary memory is not just semantic memory.  It must store structured experiment records: parent variant, child variant, mutation operator, benchmark version, random seed, model version, cost, failure mode, metric deltas, and reviewer decisions.  AutoGPT's memory can be adapted to support such records, but the project does not make them the central abstraction.  This is a general pattern across agent frameworks: memory is often designed for conversation relevance, not scientific provenance.

Tool use is one of AutoGPT's strongest contributions.  The command architecture made it easy to imagine an LLM acting through tools rather than merely answering in text.  File operations, web search, command execution, and code-related actions gave the agent a meaningful interface to the world.  From the self-evolution perspective, tools are mutation and evaluation actuators.  A self-evolving coding agent must edit code, run tests, inspect logs, compare outputs, and possibly deploy artifacts.  AutoGPT helped normalize this tool-mediated agency.  However, tool access must be governed.  Evolution systems can amplify mistakes: a faulty mutation operator combined with broad file access can corrupt its own benchmark, erase evidence, or overfit by changing tests.  Thus AutoGPT-style tools require sandboxing and policy guards.

The Forge component and Agent Protocol orientation address a related need: standardizing how agents are built and evaluated.  A protocol enables interchangeable agents and benchmark harnesses.  For self-evolution, this is essential because variants must be comparable.  If each variant changes not only its policy but also its reporting format, benchmark endpoint, or environment assumptions, the resulting comparison is invalid.  AutoGPT's movement toward platform and protocol therefore points in the right direction, even if the hype surrounding the project sometimes obscured the slower infrastructure work.

A hype-aware reading of AutoGPT should therefore separate three layers.  The first layer is the cultural layer: AutoGPT made autonomous agents visible.  The second is the architectural layer: the TAO loop, command system, memory retrieval, workspace, and later platform architecture.  The third is the evidence layer: benchmark performance, maintenance quality, contributor depth, and production reliability.  The cultural layer is unquestionably large.  The architectural layer is influential but not unique anymore.  The evidence layer is mixed and must be evaluated against current code, tests, and real deployments rather than star count.

For self-evolution, AutoGPT is best understood as a baseline substrate for autonomous exploration.  It can run iterative tasks and expose tool-mediated action.  It can host block-based workflows and persistent agents.  But an evolution layer above AutoGPT must add experiment discipline.  It should freeze a baseline agent, generate a single controlled mutation, run a benchmark suite, compare against the baseline with confidence thresholds, store the full trace, and promote only if no regression is observed.  AutoGPT's original loop says, ``think, act, observe, repeat.''  A self-evolving AutoGPT-derived platform must say, ``hypothesize, mutate, evaluate, compare, archive, and only then repeat.''

\subsection{MetaGPT: SOP-Guided Multi-Agent Collaboration}
\label{subsec:metagpt}

MetaGPT represents a different answer to the problem of agent reliability.  Instead of asking a single autonomous agent to decide everything, MetaGPT maps a software company's standard operating procedure into a multi-agent collaboration pattern.  Its slogan can be summarized as code equals SOP applied to a team.  A requirement is passed to role-specialized agents such as product manager, architect, project manager, and engineer.  Each role performs an action, emits structured artifacts, and passes messages to the next role.  The result is not just a chat transcript but a set of project documents, designs, tasks, and code artifacts.

The core abstraction is Role--Action--Message.  A role has a name, profile, goal, constraints, and available actions.  An action is a task-specific operation such as writing a product requirement document, designing an architecture, decomposing tasks, or generating code.  Messages carry the artifacts and state between roles.  This design reduces the entropy of multi-agent collaboration by assigning responsibilities.  Instead of all agents talking freely, each agent acts inside a professional frame.  In practice, this makes MetaGPT attractive for software-generation tasks because software engineering already has a rich vocabulary of roles and deliverables.

For self-evolution, MetaGPT's most valuable contribution is not merely multi-agent collaboration but \emph{procedural explicitness}.  A system can evolve only what it can represent.  If a workflow is implicit in a single long prompt, mutation is coarse and risky.  If a workflow is represented as roles, actions, message schemas, and artifact transitions, mutation can be more surgical.  One can ask whether the product-manager prompt should change, whether the architect should receive additional constraints, whether the engineer should run tests earlier, or whether a reviewer role should be inserted before code generation.  The SOP becomes an evolvable object.

The project materials note MetaGPT's academic and commercial balance.  It is an open-source framework with a large community and also has a product direction through MGX.  Its ecosystem includes research extensions such as SELA, a tree-search-enhanced approach for automated machine learning, and AFlow, an approach to automating agentic workflow generation that received ICLR oral recognition.  This matters because MetaGPT is not only a practical package but also a research platform for asking how workflows themselves can be generated or improved.  The combination of productization and academic validation gives it a stronger bridge between engineering use and self-evolution research than many purely demo-driven projects.

MetaGPT's memory system is more elaborate than a simple chat buffer.  The materials describe memory modules, brain memory, long-term memory, memory storage, and role-specific memory.  This layered memory is a natural fit for SOP-based agents because different roles require different context.  A product manager should remember requirements and user goals; an architect should remember constraints and design decisions; an engineer should remember code structure and implementation tasks.  In a self-evolving setting, role-specific memory could also store role-specific performance evidence.  For example, the architect role might accumulate examples of designs that led to test failures, while the reviewer role might accumulate bug patterns.

However, the same limitation appears again: ordinary memory is not the same as evolutionary lineage.  MetaGPT can remember messages and artifacts, but a self-evolution layer must remember variants of the SOP itself.  It must record that SOP version $v_3$ inserted an additional security review after architecture, that this change improved benchmark pass rate by two points on repository tasks but increased cost by thirty percent, and that it regressed time-to-first-code on small tasks.  Such information is not simply another message in a role's memory.  It is governance metadata about the system's own development.

MetaGPT's evaluation story is promising because its outputs are structured.  A generated project can be inspected for documents, modules, APIs, and tests.  Extensions like SELA explicitly include search, runners, evaluation, and insight modules.  AFlow addresses automatic workflow generation, which is very close to self-evolution at the orchestration level.  Nevertheless, the general-purpose MetaGPT framework still needs an external policy for persistent evaluation.  A workflow generator can produce candidate workflows, but the system must decide which candidate becomes the new default, under what evidence threshold, and with what rollback plan.

The role-based SOP model also exposes a subtle risk: institutionalizing human process assumptions that may not be optimal for agents.  Software-company roles are familiar, but LLM agents do not have the same cognitive constraints as humans.  A human product manager and architect are separated by organizational boundaries, expertise, and communication costs.  An LLM-based system may benefit from different decompositions: verifier, trace summarizer, mutation proposer, benchmark guardian, cost accountant, and regression sentinel.  Therefore, MetaGPT's insight should not be copied literally.  The deeper insight is that agent systems need explicit operating procedures.  The specific roles should be evolved and benchmarked rather than assumed.

MetaGPT is particularly relevant for \emph{team self-evolution}.  In a single-agent system, evolution may target prompts, tools, or policies.  In a team system, evolution can target division of labor.  Which roles should exist?  Which messages should they exchange?  Which artifacts should be mandatory?  Which agent has authority to reject a change?  When should a human be consulted?  These are organizational questions.  MetaGPT gives a vocabulary for answering them with code.  A self-evolution layer can then mutate that organizational code and evaluate the resulting team.

An instructive architecture is to treat MetaGPT as the execution layer for candidate SOPs.  The evolution layer stores a baseline SOP, generates a candidate change, runs both baseline and candidate on the same task suite, compares artifact quality and test pass rates, and archives all role messages.  If the candidate improves a target metric without unacceptable cost or regressions, it becomes a new SOP version.  The next generation starts from that version.  Under this design, MetaGPT's role-action-message system becomes the genotype, while project outputs and benchmark scores become the phenotype.  This genotype--phenotype framing is essential for self-evolving agent teams.

The academic and commercial balance also suggests a pragmatic lesson.  Research frameworks often optimize for novelty, while products optimize for reliability and user experience.  Self-evolution needs both.  It needs novel search over workflows, but also durable storage, user controls, observability, and deployment procedures.  MetaGPT's movement between open research, structured framework design, and productized MGX indicates the kind of ecosystem in which self-evolving systems may mature.  The remaining gap is to make evolution evidence first-class rather than an extension or experiment.

\subsection{AutoGen: Conversational Agents and Message-Driven Runtime}
\label{subsec:autogen}

AutoGen, developed by Microsoft Research, frames multi-agent work as conversation among message-driven participants.  Its newer architecture separates a core runtime from higher-level agent-chat APIs and extensions.  At the bottom is an actor-like model in which agents receive messages, process them, and emit messages through a runtime.  At the higher level are assistant agents, user proxy agents, code executor agents, group chats, round-robin teams, handoff mechanisms, and nested societies of mind.  This separation gives AutoGen both conceptual clarity and engineering flexibility.

The central idea is that an agent is a participant in a message system.  This differs from the SOP pipeline of MetaGPT and the graph-state model of LangGraph.  In AutoGen, collaboration is expressed as who says what to whom, when control shifts, and when a conversation terminates.  The runtime can support direct messages, publish-subscribe topics, subscriptions, and agent identifiers.  High-level AgentChat abstractions then package common patterns such as round-robin discussion, manager-selected group chat, or delegation to specialized agents.

This conversational orientation is powerful because many LLM capabilities are naturally dialogic.  LLMs can critique, ask clarifying questions, propose alternatives, explain reasoning, and revise answers in response to peer messages.  AutoGen turns this into a programming model.  A writer agent can draft, a reviewer agent can critique, a coder agent can execute, and a user-proxy agent can represent human approval.  A Society of Mind agent can encapsulate an internal team behind a single external interface.  Such nesting is particularly useful for complex tasks because it allows a system to expose a simple surface while hiding internal deliberation.

For self-evolution, AutoGen's message-driven architecture is valuable in three ways.  First, it makes deliberation traceable.  Variant proposals, critiques, votes, and handoffs can be represented as messages.  Second, it supports multiple coordination policies.  A self-evolution system can compare round-robin review against manager-selected review or specialist handoff.  Third, the actor-like runtime can in principle distribute agents across processes or languages, which matters when an evolution pipeline includes heterogeneous evaluators, code runners, and monitoring services.

AutoGen's tool system also maps well to evolutionary workflows.  FunctionTool wraps callable capabilities.  Workbench manages collections of tools.  Tool agents can separate tool execution from ordinary conversation.  CodeExecutorAgent supports code-related tasks.  In a self-evolving system, this separation can enforce safety.  The mutation proposer need not execute arbitrary shell commands; it can request execution through a controlled tool agent.  The benchmark agent can be the only participant allowed to write benchmark results.  The promotion agent can be the only participant allowed to update a baseline.  AutoGen's handoff and team abstractions are thus not merely conveniences; they can encode authority boundaries.

The memory and model-context abstractions are lighter than a full experiment database, but they provide hooks.  ListMemory and model contexts can store conversation facts.  More importantly, AutoGen's messages can be logged externally as structured traces.  A self-evolution layer should treat AutoGen conversations as evidence streams.  Every proposal, critique, tool call, result, and termination reason should be stored with run identifiers.  Later analysis can ask which reviewer comments predicted real failures, which handoff policies reduced cost, and which agents contributed to successful variants.  This turns conversation from ephemeral coordination into analyzable evolutionary data.

AutoGen's evaluation support includes agbench and task-specific evaluation patterns.  Its Studio interface also helps users inspect and build workflows.  Yet, as with the other frameworks, evaluation is not the same as evolution.  AutoGen can run a group chat that improves an answer, but persistent self-evolution requires comparing group-chat configurations across tasks.  For example, does adding a critic agent improve accuracy enough to justify cost?  Does a manager-selected speaker policy outperform round-robin on debugging tasks?  Does a Society of Mind wrapper hide useful trace details from the outer evaluator?  These are empirical questions that require systematic experiment design.

The framework is especially suitable for \emph{social search} over agent configurations.  In social search, candidate systems differ not only in prompts but also in interaction topology.  One candidate may use a proposer--critic--executor trio.  Another may use two independent solvers followed by an adjudicator.  A third may use a manager that dynamically selects specialists.  AutoGen can express all three as conversational teams.  An evolution layer can then mutate team composition, speaker policy, termination conditions, memory access, and tool workbenches.  Fitness can be measured by benchmark success, cost, latency, and human approval.

A limitation of conversational frameworks is that conversation can become an attractor in itself.  Agents may discuss without acting, critique without testing, or agree prematurely.  This is not a flaw unique to AutoGen; it is a property of language-agent systems.  For self-evolution, conversation must be tied to external evidence.  A reviewer message should not count as proof unless the reviewer cites a test, trace, or formal check.  A claimed improvement should trigger a benchmark run.  A handoff should be logged with the reason for transfer and the outcome.  Without such grounding, a conversational team can evolve rhetoric rather than capability.

AutoGen's layered design offers a useful template for self-evolution architecture.  The core layer provides message transport and agent lifecycle.  The chat layer provides convenient teams.  The extension layer provides tools and model integrations.  A self-evolution layer can be placed above these: it defines experiment plans, assigns agents to propose and evaluate variants, stores traces, and promotes configurations.  This preserves AutoGen's strengths while adding the missing governance layer.  The result is not an AutoGen replacement but an AutoGen-based evolutionary laboratory.

\subsection{CrewAI: Lightweight Crews, Production Flows, and Community Health}
\label{subsec:crewai}

CrewAI takes a pragmatic position in the agent-framework landscape.  It is designed as an independent, lightweight, high-performance framework rather than a thin layer over a larger agent stack.  The source materials emphasize that it is built without depending on LangChain or another agent framework.  This zero-dependency stance is not merely branding.  It affects how developers reason about the system: fewer inherited abstractions, fewer transitive dependencies, and more direct control over the agent, task, and process model.

The core API is organized around agents, tasks, and crews.  An agent has a role, goal, backstory, tools, and optional memory.  A task has a description, expected output, and assigned agent.  A crew combines agents and tasks under a process such as sequential execution or hierarchical management.  This model is easy to teach and easy to inspect.  It maps naturally onto business workflows where responsibilities are known in advance.  A research agent gathers information, a writer agent drafts, a reviewer agent critiques, and a manager agent can coordinate a hierarchy.

CrewAI's most distinctive architectural point is the split between Crews and Flows.  Crews represent autonomous collaboration among agents.  Flows represent event-driven, production-grade orchestration where the developer controls transitions more precisely.  This dual architecture acknowledges a tension that appears throughout the field.  Autonomous agents are flexible but hard to predict.  Deterministic workflows are predictable but less adaptive.  Production systems often need both: autonomous subroutines embedded in a controlled outer flow.  CrewAI makes this split explicit.

For self-evolution, the Crew + Flow distinction is highly useful.  The evolution layer itself should usually be a controlled flow, because mutation, evaluation, comparison, and promotion must occur in a disciplined order.  Inside that flow, crews can be used to solve tasks, critique variants, or generate hypotheses.  For example, a self-evolution flow might start with benchmark selection, then invoke a proposal crew, then run candidate agents, then call a reviewer crew, then execute statistical comparison, and finally decide promotion.  The flow prevents the system from skipping evaluation, while the crew provides flexible reasoning within each stage.

CrewAI's memory model includes short-term, long-term, and entity memory, plus knowledge management.  This is richer than a single conversation buffer and aligns with the needs of repeated automation.  Entity memory is especially relevant for long-running tasks: agents may need to remember users, repositories, APIs, datasets, or organizations.  Long-term memory can preserve lessons across tasks.  However, as in other frameworks, evolutionary memory requires additional structure.  The system must not merely remember that a tool worked before; it must remember under which variant, benchmark, and environment that tool contributed to improvement.

The project materials describe CrewAI as having some of the healthiest community metrics among the surveyed frameworks.  This claim should be interpreted qualitatively rather than as a single-number ranking.  Compared with a hype-heavy project whose stars vastly outpace contributor depth, CrewAI's signals include a focused framework identity, substantial developer education, a cloud/control-plane story, and an API designed for repeated production use.  Its community health appears tied to actual workflow building rather than only viral curiosity.  For self-evolution research, this matters because a framework's practical viability depends on documentation, examples, maintenance, and user trust.

The zero-dependency design also has implications for evolvability.  A framework with fewer external conceptual dependencies may be easier to mutate and reason about.  If an evolved workflow fails, the cause is less likely to be hidden in a deep stack of inherited abstractions.  On the other hand, zero dependency can mean reimplementing capabilities that larger ecosystems already provide.  The self-evolution designer must decide whether control or ecosystem breadth is more important.  CrewAI favors control and simplicity; LangGraph favors integration with LangChain's ecosystem; AutoGen favors runtime and conversation architecture; DSPy favors optimizer abstractions.

CrewAI's evaluation layer is less central than DSPy's.  Tasks specify expected outputs, and control-plane observability can help monitor runs, but the framework does not by itself define a scientific evaluation loop.  Therefore, a CrewAI-based self-evolution system should attach explicit evaluators to flows.  Each task should produce machine-checkable outputs where possible.  Agents should not decide their own success unless their decision is itself evaluated.  Flow steps should record metrics, costs, and traces.  Promotion should occur only after comparison with a frozen baseline crew or flow.

A useful mutation space for CrewAI includes agent role descriptions, goals, backstories, tool assignments, process type, task ordering, manager policies, and flow transitions.  Some mutations are local: changing a researcher's tool set.  Others are structural: switching from sequential to hierarchical execution.  Still others are governance-oriented: adding a reviewer task before final output.  Because CrewAI's concepts are high-level and human-readable, candidate mutations can often be explained clearly.  This improves human oversight.  A reviewer can understand that the evolved system added a fact-checking agent or moved testing earlier in the flow.

CrewAI also illustrates an important engineering principle: self-evolution should be accessible to ordinary developers.  If evolving an agent team requires deep knowledge of distributed runtimes, graph algorithms, and optimizer internals, adoption will be limited.  CrewAI's simple crew/task vocabulary can lower the barrier.  A self-evolution layer built on top of CrewAI could present improvements as changes to roles, tasks, and flows, rather than opaque parameter vectors.  This does not eliminate the need for rigorous evaluation, but it makes the evolved artifact easier to audit.

The main risk is that intuitive abstractions can hide quantitative uncertainty.  A new role may sound useful but degrade performance.  A hierarchical manager may improve coordination on large tasks but slow down small tasks.  Long-term memory may help personalization but introduce stale assumptions.  Therefore, CrewAI's production-friendly design should be paired with DSPy-like metrics and LangGraph-like checkpoints.  In a combined architecture, CrewAI supplies the work organization, LangGraph supplies the outer evolutionary flow, DSPy optimizes prompts inside roles, and an external ledger supplies lineage and regression control.

\subsection{DSPy: Declarative LLM Programming and Compiled Optimization}
\label{subsec:dspy}

DSPy is the most directly optimization-oriented framework in this survey.  Its name, Declarative Self-improving Python, captures its central premise: LLM applications should be written as declarative programs whose prompts and demonstrations can be compiled against metrics.  Instead of hand-crafting a prompt string and hoping it generalizes, the developer defines signatures, modules, and optimizers.  The framework then searches for better instructions, demonstrations, or program configurations using training and development examples.

The basic abstraction stack is Signature--Module--Optimizer.  A signature declares inputs and outputs, such as question to answer or document to summary and citation.  A module composes predictors such as ChainOfThought, ReAct, ProgramOfThought, Refine, Retry, or BestOfN.  An optimizer, historically called a teleprompter, compiles a module against data and a metric.  This separation between declaration and execution is a major conceptual contribution.  The developer specifies what the program should do and how it is composed, while the optimizer searches for how prompts and demonstrations should be instantiated.

For self-evolution, DSPy's key insight is that prompts are program parameters.  In much of agent engineering, prompts are treated as text artifacts edited manually.  DSPy treats them more like weights, examples, or compiler outputs.  They can be generated, scored, selected, and saved.  This makes improvement more systematic.  If a metric is available, the system can optimize rather than merely reflect.  The difference is profound: reflection produces plausible advice, whereas optimization requires measured improvement.

DSPy's optimizer family provides several evolutionary analogues.  BootstrapFewShot uses a teacher model and training examples to collect successful traces that become demonstrations.  MIPROv2 jointly searches instructions and demonstrations, using an optimization backend such as Bayesian search.  COPRO explores instruction candidates.  SIMBA uses stochastic minibatches, high-variance examples, and self-reflection to generate improvement rules or add successful examples.  BootstrapFinetune collects trajectories for model fine-tuning.  GRPO introduces reinforcement-learning-style optimization.  These methods are not identical to genetic algorithms, but they implement the central evolutionary pattern: generate candidates, evaluate against a metric, and preserve better configurations.

DSPy's Evaluate abstraction is equally important.  It makes metrics first-class.  A program is not improved because an agent says it is improved; it is improved because it scores better on a development set under a defined metric.  The evaluator can run examples in parallel, apply thresholds, and return structured results.  This is closer to scientific self-evolution than most agent frameworks.  The system has an explicit objective function, a compilation procedure, and artifacts that can be compared.

However, DSPy is not a full autonomous-agent infrastructure layer.  It does not primarily solve long-running tool orchestration, sandboxed software development, human approvals, or multi-agent social coordination.  It can include ReAct-style modules and tool-using predictors, but its center of gravity is program optimization.  This means DSPy is best seen as an optimization engine inside a broader self-evolution system.  It can optimize prompts and modules used by agents in AutoGen, CrewAI, MetaGPT, or LangGraph.  The broader system can handle task scheduling, environment execution, trace storage, and deployment.

The separation between declaration and execution is especially valuable for reproducibility.  If a program is written declaratively, the optimizer can compile different versions while preserving a stable interface.  The input and output schema remains fixed, while instructions and demonstrations change.  This allows controlled experiments.  A self-evolution system can say: the signature and metric are unchanged; only the instruction set changed.  That makes attribution clearer.  By contrast, if a free-form agent prompt changes its goals, output format, and evaluation behavior simultaneously, improvement is hard to interpret.

DSPy also encourages a healthy humility about prompt engineering.  Human prompt authors often overfit to a few examples and infer broad principles from anecdote.  DSPy's compilation loop forces prompts to face data.  In self-evolution terms, this is the difference between Lamarckian storytelling and empirical selection.  The system may generate a rule explaining why a prompt should improve, but that rule matters only if it improves measured performance.  This evidence discipline should be imported into agent frameworks more broadly.

The limitations are also instructive.  DSPy's optimization quality depends on metrics and datasets.  If the metric is weak, the optimizer can overfit or optimize the wrong behavior.  If the development set is narrow, the compiled prompt may not generalize.  If the task requires open-ended software environment interaction, a simple input-output metric may miss important process failures.  Therefore, a self-evolution platform should treat DSPy compilation as one stage in a larger evaluation funnel.  A prompt that improves unit examples should still be tested inside the real agent workflow, under real tool conditions, with regression checks.

Another challenge is lineage.  DSPy produces optimized programs and can record evaluation results, but a full self-evolution system needs a durable artifact registry.  Each compiled program should be assigned an identifier, linked to parent programs, associated with optimizer settings, training data versions, metric versions, model versions, and cost.  Without this registry, successful optimizations become difficult to reproduce.  With it, DSPy can serve as the compiler inside an evolutionary software supply chain.

A useful architecture is to embed DSPy at the leaves of agent systems.  In a MetaGPT team, each role's action prompt can be a DSPy module.  In a CrewAI crew, each agent's task-specific reasoning can be compiled.  In AutoGen, specialized assistant agents can use DSPy-optimized modules for particular message types.  In LangGraph, nodes can wrap DSPy programs and feed their outputs into graph state.  In OpenHands, code-review or bug-localization prompts can be optimized against repository benchmarks.  The evolution layer then coordinates when to compile, how to evaluate compiled artifacts, and when to promote them.

DSPy therefore represents the clearest path from prompt engineering to prompt evolution.  Its lesson for the broader field is that self-improvement must be operationalized as compilation against evidence.  Natural-language reflection remains useful for generating candidates and interpreting failures, but it must be subordinate to metrics.  In the architecture of AI self-evolution, DSPy is not the whole organism; it is a specialized organ for optimizing LLM program components.

\subsection{LangGraph: Graph-Based Agent Workflows and Cyclic Refinement}
\label{subsec:langgraph}

LangGraph approaches agent infrastructure through graph execution.  A workflow is modeled as a directed graph whose nodes are functions or agents and whose edges define control flow.  State is carried through channels.  Conditional edges route execution based on the current state.  Cycles are allowed.  Checkpoints can persist state so that execution can pause, resume, or recover.  This makes LangGraph one of the most natural substrates for iterative refinement workflows.

The central abstraction is StateGraph.  A developer defines a state schema, adds nodes, adds fixed or conditional edges, and compiles the graph into an executable object.  Under the hood, LangGraph uses a Pregel-inspired execution model with supersteps, channels, and checkpointing.  The framework supports state persistence through stores such as SQLite or PostgreSQL, and it integrates with LangChain tools and prebuilt agents.  Studio and Cloud offerings support visualization and deployment.

For self-evolution, graph structure is extremely attractive because evolution workflows are not simple chains.  They contain loops, branches, retries, stopping conditions, and human gates.  A candidate may fail syntax checks and return to a repair node.  A benchmark may reveal a regression and route to rollback.  A high-cost improvement may require human approval.  A flaky result may trigger reruns.  A promising variant may branch into further local search.  Encoding these behaviors as an explicit graph is clearer than hiding them inside a monolithic agent prompt.

LangGraph's support for cycles is especially important.  Many workflow systems assume acyclic DAGs, which are useful for data pipelines but awkward for refinement.  Self-evolution is inherently cyclic: propose, test, diagnose, repair, retest.  Cycles must be controlled, however.  The graph should include iteration limits, convergence checks, cost budgets, and failure exits.  LangGraph provides the structural capacity for cycles; the evolution designer must provide the governance rules.

Checkpointing is another major contribution.  In ordinary agent runs, failure often destroys context or forces the system to repeat expensive steps.  In an evolutionary pipeline, this is unacceptable.  If a candidate passes early tests and fails later tests, the system should retain its state for diagnosis.  If a human approval is needed, the run should pause without losing trace data.  If infrastructure fails, the pipeline should resume from a checkpoint.  LangGraph's checkpoint model directly supports these needs.

State channels also encourage disciplined data flow.  Instead of relying on a single chat transcript, a graph can maintain separate fields for task input, candidate patch, benchmark results, reviewer notes, cost, risk flags, and promotion decision.  This is a better fit for self-evolution than undifferentiated text memory.  Structured state makes it easier to enforce policies: the promotion node should read benchmark results and regression flags, not merely a natural-language summary.  It also makes audit easier because the path through the graph records why a decision occurred.

LangGraph does not, by itself, decide what constitutes improvement.  It provides orchestration and persistence, not an evolutionary objective.  Metrics, datasets, mutation operators, and selection rules must be supplied externally.  This is not a weakness; it is a clean separation of concerns.  LangGraph can be the workflow engine for an evolution controller, while DSPy supplies prompt compilation, OpenHands supplies sandboxed code execution, and AutoGen or CrewAI supplies deliberative teams.  In this combined architecture, LangGraph is the spine.

A typical self-evolution graph might contain nodes such as baseline freeze, mutation proposal, candidate materialization, static validation, benchmark execution, trace summarization, regression comparison, cost analysis, human review, promotion, and archival.  Conditional edges would route syntax failures back to repair, benchmark failures to diagnosis, inconclusive results to rerun, and successful results to promotion.  Checkpoints would be saved after every expensive stage.  The graph's state would include lineage identifiers and artifact paths.  Such a graph makes the evolution process inspectable and repeatable.

The risk of graph-based systems is overengineering.  A graph can become complex, with many nodes and conditional routes that are themselves hard to understand.  Evolution can then target the graph and accidentally create unmaintainable control flow.  To prevent this, graph mutations should be constrained.  The system might allow adding a reviewer node or changing a routing threshold, but not arbitrary rewiring without validation.  Graph schemas and invariants are essential.  For example, every path to promotion must pass through benchmark comparison; every mutation must have a parent; every failed run must be archived.

LangGraph also changes how we think about agent autonomy.  In a loop-centered agent, autonomy lives inside the model's action selection.  In a graph-centered system, autonomy can be localized to particular nodes while the outer structure remains deterministic.  This is often preferable for self-evolution.  The mutation proposer can be creative; the benchmark runner should be deterministic; the promotion policy should be conservative.  Graphs allow these different epistemic modes to coexist.

Because LangGraph is part of the LangChain ecosystem, it benefits from a broad tool and model integration base.  This is valuable but can introduce dependency complexity.  A self-evolution system using LangGraph should record versions of tools, model clients, checkpoint stores, and graph definitions.  Otherwise, improvements may be confounded by changing dependencies.  The more powerful the ecosystem, the more important provenance becomes.

LangGraph's highest value for self-evolution is therefore not that it is an agent framework in the conversational sense, but that it is an execution grammar for cyclic, stateful, inspectable improvement.  It can express the difference between trying again and promoting a new baseline.  It can require that evidence pass through structured channels.  It can make human approval a node rather than an afterthought.  These properties are central to moving self-evolving AI from demos to engineering practice.

\subsubsection*{Complementary case: CAMEL role-playing societies, critics, and environment-rich agents}

CAMEL is one of the earliest frameworks to explore communicative agents through role-playing.  Its original research direction asked what happens when LLMs are assigned roles and asked to collaborate through structured dialogue.  The framework has since expanded into a broad agent ecosystem with ChatAgent, RolePlaying societies, Workforce orchestration, critic agents, toolkits, interpreters, retrievers, storages, environments, verifiers, runtimes, and data generation utilities.  This breadth makes CAMEL particularly relevant for research on agent societies and self-improvement.

The RolePlaying abstraction is the signature pattern.  A task can be refined by a task-specification agent, planned by a task-planning agent, and then pursued by a user-role agent and assistant-role agent in dialogue.  A critic can be placed in the loop to evaluate or improve the interaction.  This structure is less rigid than MetaGPT's software-company SOP but more structured than free-form chat.  It provides a controlled social environment in which roles shape behavior.

For self-evolution, the critic-in-the-loop pattern is highly relevant.  Many self-improvement methods begin with an agent reflecting on its own output.  CAMEL externalizes part of that reflection into a critic role.  This is often better than self-critique by the same agent, because the critic can have a different prompt, memory, tool access, or evaluation standard.  A critic can identify failure modes, request revisions, and provide training data for future improvements.  However, critics also need evaluation.  A persuasive critic is not necessarily a correct critic.  The evolution layer must measure whether critic involvement improves downstream metrics.

CAMEL's memory system includes chat-history memory, score-based context creation, and memory records with metadata.  Score-based context selection is useful for long conversations where not every prior message should be included.  In self-evolving systems, context selection itself can become an optimization target.  Which memories should be retrieved for a debugging task?  Which prior failures should inform a planning agent?  Which critic comments should become durable rules?  CAMEL provides primitives for memory selection, but an evolution layer can add empirical comparison among memory policies.

The environment and tool subsystems are broad.  CAMEL includes code interpreters such as subprocess, Docker, Jupyter, and microsandbox variants; toolkits for OpenAPI, terminal, and browser operations; retrievers including vector, BM25, and hybrid retrieval; storage backends including vector, graph, key-value, and object stores; and environments such as single-step and multi-step settings.  This makes CAMEL suitable for experiments that require more than text generation.  A self-evolving agent can act, verify, retrieve, and store evidence across different environment types.

The Self-Instruct data generation component is another connection to self-evolution.  Generating instruction data is a form of improving future training or tuning resources.  An agent society can generate tasks, solve them, critique them, and filter them.  Such synthetic-data loops can be powerful but dangerous.  Without external grounding, they may amplify biases or hallucinations.  Therefore, CAMEL-style data generation should be coupled with verifiers, benchmarks, and diversity controls.  The generated data should have lineage just like evolved prompts or workflows.

CAMEL's Workforce abstraction provides a more general multi-agent organization mechanism.  Agents can be added to a workforce and assigned tasks.  This overlaps with CrewAI and AutoGen but retains CAMEL's research orientation and broad environment support.  A self-evolution system can use Workforce to assemble candidate teams, then evaluate them under controlled tasks.  Candidate teams might differ in role assignment, critic placement, planner use, or tool distribution.  The evolution layer can select among these social configurations.

One of CAMEL's strengths is that it treats agent societies as experimental objects.  Role names, system messages, task-specification settings, planner settings, critic settings, and environment configurations are all variables.  This aligns closely with the idea of evolving agents.  The challenge is to avoid combinatorial explosion.  The number of possible role and tool configurations is enormous.  A self-evolution layer needs search-space management: define mutation operators, restrict unsafe combinations, use cheap proxy metrics before expensive benchmarks, and preserve diversity without losing focus.

The verifier modules in CAMEL point toward an important principle: language-based evaluation should be supplemented with formal or executable checks where possible.  A math verifier or Python verifier can ground an agent's output.  In self-evolution, such verifiers can serve as fitness functions.  For example, a generated solution either passes a Python test or does not.  A critic's natural-language judgment may help diagnose failure, but the verifier provides selection pressure.  This mirrors the broader argument of the chapter: self-evolution requires evidence, not just reflection.

CAMEL also highlights the distinction between research flexibility and production governance.  A research framework benefits from many agent types, tools, storages, and environments.  A production self-evolution system must constrain this richness.  It must know which tools are allowed during mutation, which storages are authoritative, which environments are reproducible, and which generated artifacts can be promoted.  CAMEL provides a rich laboratory; the evolution layer must provide lab protocol.

In a combined self-evolution architecture, CAMEL can serve as the social experimentation engine.  It can generate role-playing dialogues, critic feedback, synthetic tasks, and candidate team behaviors.  LangGraph can orchestrate the outer loop.  DSPy can optimize specific prompts.  OpenHands can execute software tasks.  An experiment ledger can store lineage.  Under this composition, CAMEL contributes diversity and critique, while other layers enforce measurement and safety.

\subsection{Evolution Capability Gaps}
\label{subsec:evolution-capability-gaps}

Across the surveyed frameworks, the most important finding is a gap between \emph{agent capability} and \emph{evolution capability}.  Agent capability means that a framework can run agents, call tools, manage memory, coordinate multiple participants, or execute workflows.  Evolution capability means that a system can produce successive variants, evaluate them under controlled conditions, preserve lineage, prevent regressions, and promote improvements.  Most current frameworks are strong on the former and partial on the latter.

\begin{table*}[t]
\centering
\small
\caption{Common evolution capability gaps in current agent frameworks.}
\label{tab:evolution-gaps}
\begin{tabular}{p{0.18\textwidth}p{0.30\textwidth}p{0.24\textwidth}p{0.20\textwidth}}
\hline
\textbf{Gap} & \textbf{Typical current behavior} & \textbf{Why it matters for self-evolution} & \textbf{Needed layer} \\
\hline
Persistent evaluation & Evaluation is task-specific, ad hoc, or external; some frameworks provide benchmarks or metrics but not universal governance. & A variant cannot be trusted without repeated, comparable measurement. & Benchmark registry, evaluator API, confidence and rerun policy. \\
Lineage tracking & Runs may log traces, but parent--child variant relationships are not always first-class. & Improvement requires knowing what changed and what evidence justified promotion. & Artifact registry with parent IDs, mutation metadata, and versioned configs. \\
Regression prevention & Agents may improve one task while silently harming another. & Evolution without regression checks causes capability drift and benchmark overfitting. & Baseline freeze, regression suite, rollback rules, gated promotion. \\
Cross-run memory & Memory often supports a task or conversation rather than long-term experiment science. & Lessons must survive beyond a single run and remain attributable. & Structured experiment memory, failure taxonomy, learned rule database. \\
Cost and risk accounting & Frameworks expose token or runtime information unevenly. & A small quality gain may be unacceptable if cost, latency, or risk explodes. & Multi-objective fitness with cost, latency, safety, and maintainability. \\
Reproducible environments & Tool and sandbox behavior may vary across runs. & Variant comparisons are invalid if environments drift. & Environment snapshots, seed control, dependency manifests. \\
Promotion governance & A successful output may be accepted manually or implicitly. & Self-evolution needs a policy for when a child becomes the new baseline. & Promotion controller, approval workflow, audit log. \\
\hline
\end{tabular}
\end{table*}

The first gap is persistent evaluation.  DSPy has the strongest built-in evaluation loop, and AutoGen includes agbench-related tooling.  Other frameworks often rely on user-defined expected outputs, tests, or external benchmarks.  But self-evolution requires evaluation to be persistent and comparable across generations.  The system must know which benchmark version was used, which examples were included, which metric was applied, and whether the result was statistically stable.  A one-off successful run is not evolution; it is an anecdote.

The second gap is lineage tracking.  Current frameworks can often save runs, messages, checkpoints, or artifacts.  They less often treat the framework configuration itself as an evolving artifact with parent and child relationships.  A lineage record should answer: what was the parent? what mutation was applied? who or what proposed it? what files, prompts, roles, graph edges, tools, or demonstrations changed? what evidence was collected? why was it accepted or rejected? Without this record, the system cannot learn which mutation operators are productive.

The third gap is regression prevention.  Agent systems are prone to local improvements.  A prompt change may improve one benchmark while harming another.  A memory policy may help long tasks while confusing short tasks.  A new critic may catch more bugs but also introduce unnecessary revisions.  A self-evolution platform must freeze safe baselines and run regression suites.  It should treat promotion as a high-risk operation, not an automatic consequence of one improved score.  The conservative rule should be: no promotion without baseline comparison and regression evidence.

The fourth gap is cross-run memory.  Framework memory is often designed to answer the question, ``What does this agent need to know now?'' Evolution memory must answer, ``What has the system learned from many experiments?'' These are different.  Cross-run memory should store failure modes, mutation outcomes, benchmark sensitivities, environment issues, and reviewer reliability.  It should be queryable by future mutation proposers.  It should also distinguish durable lessons from stale accidents.  A failed run caused by a temporary API outage should not become a permanent rule against a method.

The fifth gap is multi-objective accounting.  Many agent examples optimize task success or output quality.  Production self-evolution must optimize quality under constraints.  Cost, latency, token usage, tool risk, security exposure, maintainability, and human-review burden all matter.  A candidate that improves accuracy by one percent while doubling cost may or may not be acceptable depending on the use case.  Therefore, the fitness function should be explicit and multi-dimensional.  Frameworks expose some telemetry, but few make multi-objective promotion a core abstraction.

The sixth gap is environment reproducibility.  Tool-using agents interact with changing worlds: websites change, package versions shift, model APIs update, file systems differ, and random seeds alter outputs.  Evolutionary comparison is invalid when the environment is uncontrolled.  OpenHands-style sandboxes, LangGraph checkpoints, and AutoGPT platform runtimes help, but the evolution layer must record environment manifests.  Ideally, candidate variants run in isolated, reproducible sandboxes with controlled dependencies and model versions.

The seventh gap is promotion governance.  In many frameworks, success is local: the agent returns an answer, the workflow ends, or the user accepts an output.  In self-evolution, success can change the future system.  This requires governance.  Who is allowed to promote a new prompt?  Can an agent update its own evaluator?  Must a human approve changes to tools?  What happens if a promoted variant later regresses?  These questions are outside the normal scope of agent orchestration, but they are central to safe self-evolution.

These gaps explain why existing frameworks should be viewed as components rather than complete self-evolving systems.  They provide execution substrates.  They do not provide the full evolution operating system.  A self-evolution platform needs a separate control plane containing a variant registry, evaluation service, trace store, memory system, regression suite, promotion policy, and rollback mechanism.  The agent framework is then plugged into this control plane as a runner.

\subsection{Self-Evolve Positioning: An Evolution Layer Above Frameworks}
\label{subsec:self-evolve-positioning}

The appropriate positioning for a Self-Evolve platform is above existing frameworks, not beside them as yet another agent framework.  The ecosystem already contains many ways to define agents, graphs, crews, conversations, tools, and software environments.  Rebuilding all of them would fragment effort.  The missing layer is the one that turns these substrates into evidence-driven evolutionary systems.  Self-Evolve should therefore act as an evolution control plane.

Such a control plane begins with a uniform representation of variants.  A variant might be a DSPy compiled module, a LangGraph graph definition, a CrewAI flow, an AutoGen team, a MetaGPT SOP, a CAMEL role-playing configuration, an AutoGPT block workflow, or an OpenHands microagent.  The representation must capture the framework type, configuration, artifact paths, dependency versions, model choices, tool permissions, and parent lineage.  The exact internal structure can remain framework-specific, but the outer metadata should be uniform.

The next component is a mutation layer.  Mutations should be typed.  A prompt mutation changes instructions or demonstrations.  A topology mutation changes roles, graph nodes, or team members.  A tool mutation adds, removes, or restricts tools.  A memory mutation changes retrieval policy or persistence.  A workflow mutation changes ordering or routing.  A code mutation changes implementation.  Typed mutations make evaluation interpretable.  They also allow safety policies: for example, tool-permission mutations may require human approval, while demonstration mutations may be allowed automatically.

The third component is an evaluator layer.  It should support unit-style tests, benchmark suites, human review, model-based grading, formal verifiers, cost analysis, and adversarial checks.  Evaluation should be repeatable and versioned.  Every score should be linked to the evaluator version and dataset version.  The evaluator should compare candidates against a frozen baseline, not merely against an absolute threshold.  It should produce both aggregate metrics and per-example traces so that later agents can diagnose failures.

The fourth component is an experiment ledger.  This ledger is the memory of evolution.  It stores variants, mutations, runs, traces, metrics, costs, failures, decisions, and promotions.  It should be queryable by future agents.  For example, before proposing a mutation, an agent should be able to ask which previous mutations helped similar tasks, which failed due to cost, and which caused regressions.  This prevents the system from rediscovering the same failed ideas.  It also enables meta-evolution: improving the mutation strategy itself based on historical yield.

The fifth component is a promotion and rollback policy.  Promotion should be gated.  A candidate might need to beat the baseline by a minimum delta, pass all critical regression tests, remain within cost limits, and receive approval for high-risk changes.  Rollback should be automatic when post-promotion monitoring detects regression.  The platform should keep enough lineage information to revert not only files but also prompts, graph definitions, memory policies, and tool permissions.

The sixth component is an interface to existing frameworks.  Rather than forcing all users into one abstraction, Self-Evolve should provide adapters.  A DSPy adapter compiles modules and reports metrics.  A LangGraph adapter runs graphs with checkpoints.  A CrewAI adapter executes crews and flows.  An AutoGen adapter captures messages and team configurations.  A MetaGPT adapter treats SOPs as evolvable artifacts.  A CAMEL adapter runs role-playing societies and extracts critic evidence.  An OpenHands adapter executes code tasks in sandboxes.  An AutoGPT adapter runs autonomous or block-based agents.  The evolution layer normalizes evidence across them.

\begin{table*}[t]
\centering
\small
\caption{Recommended positioning of Self-Evolve relative to existing frameworks.}
\label{tab:self-evolve-positioning}
\begin{tabular}{p{0.20\textwidth}p{0.34\textwidth}p{0.34\textwidth}}
\hline
\textbf{Layer} & \textbf{Responsibility} & \textbf{Examples} \\
\hline
Framework runtime & Execute agents, workflows, tools, conversations, or code environments. & AutoGPT loop, MetaGPT SOP, AutoGen team, CrewAI crew/flow, DSPy module, CAMEL society, LangGraph graph, OpenHands sandbox. \\
Evolution adapter & Translate framework-specific configs, traces, and outputs into a common experiment schema. & Extract prompts, roles, graph edges, tool lists, checkpoints, messages, costs, and artifacts. \\
Mutation engine & Generate typed candidate changes under safety constraints. & Prompt edits, role additions, graph-route changes, memory-policy changes, tool restrictions, code patches. \\
Evaluation engine & Run candidates against versioned tasks and compare with baselines. & Unit tests, benchmark suites, LLM judges, verifiers, human review, cost and latency metrics. \\
Lineage ledger & Preserve parent--child relations, traces, metrics, and decisions. & Variant registry, run records, failure taxonomy, promotion history, rollback pointers. \\
Governance policy & Decide whether a candidate can become a new baseline. & Regression gates, confidence thresholds, budget limits, approval rules, rollback triggers. \\
\hline
\end{tabular}
\end{table*}

This layered positioning also clarifies the relationship between self-evolution and AutoML.  AutoML systems search over model architectures, hyperparameters, preprocessing, or pipelines.  A Self-Evolve system searches over agentic programs: prompts, tools, memories, roles, graphs, workflows, and code.  The search space is more heterogeneous and more safety-sensitive.  Therefore, Self-Evolve needs AutoML-like rigor but agent-specific abstractions.  It must understand conversations, tool permissions, memory retrieval, and workflow topology, not only numeric hyperparameters.

The layer should also support different time scales.  Short-horizon adaptation happens within a run: retry, refine, ask a critic, repair a patch.  Medium-horizon optimization happens across a batch of examples: compile a prompt, select demonstrations, tune a graph threshold.  Long-horizon evolution happens across releases: promote a new agent team, update a workflow, retire a tool, or change the default memory policy.  Current frameworks often support the first time scale and sometimes the second.  Self-Evolve should unify all three.

A practical Self-Evolve pipeline might proceed as follows.  First, freeze a baseline configuration and its environment.  Second, select a mutation dimension, such as a role prompt or graph edge.  Third, generate a candidate with a bounded mutation operator.  Fourth, materialize the candidate in the target framework.  Fifth, run a smoke test to catch invalid configurations.  Sixth, run a benchmark suite against both baseline and candidate.  Seventh, compare results with statistical and cost thresholds.  Eighth, archive traces and failure cases.  Ninth, promote, reject, or send the candidate to repair.  Tenth, update cross-run memory with the lesson learned.  This process is independent of whether the underlying executor is CrewAI, LangGraph, AutoGen, or another framework.

The survey also suggests that no single framework should dominate the evolution layer.  DSPy is strongest for prompt/module optimization.  LangGraph is strongest for explicit cyclic orchestration.  OpenHands is strongest for real software-environment actions.  AutoGen is strong for message-driven multi-agent deliberation.  CrewAI is strong for accessible production workflows.  MetaGPT is strong for SOP-based team structure.  CAMEL is strong for role-playing societies and critics.  AutoGPT is historically strong for autonomous loop exploration and platform visibility.  A mature Self-Evolve platform should be pluralistic and adapter-based.

Finally, Self-Evolve should treat human oversight as part of the architecture, not as an emergency patch.  Some mutations are low risk and can be promoted automatically after tests.  Others change tools, permissions, external integrations, or evaluation criteria.  These should require human review.  Frameworks such as LangGraph already make human-in-the-loop nodes natural; AutoGen and CrewAI can represent review agents or user proxies; OpenHands can show code diffs and terminal traces.  The evolution layer should standardize when and how human approval is required.

In conclusion, agent frameworks have reached a level of maturity where they can execute impressive tasks, coordinate multiple agents, and integrate powerful tools.  The next bottleneck is not merely more autonomy.  It is disciplined improvement.  The field needs systems that remember what they tried, measure what changed, prevent regressions, and govern promotion.  AutoGPT, MetaGPT, AutoGen, CrewAI, DSPy, CAMEL, LangGraph, and OpenHands each provide essential pieces of this future.  The Self-Evolve layer is the missing connective tissue: an evidence-first evolutionary control plane that turns agent execution into cumulative, auditable, and safe capability growth.
